% Book 3, Year 1, chapter 28 -- sequential circuits and finite-state machines.
% Provenance (never printed): sources/book3/BRIEFS.md, chapter 28.
% ledger: b1:mealy-1955, b1:moore-1956, b1:core-memory
\chapter{Sequential Circuits and Finite-State Machines}\label{ch:b1:sequential-circuits}

A traffic light changes every thirty seconds, and it knows which colour comes next. The circuits of
\cref{ch:b1:combinational-circuits} cannot do that: their outputs depend only on their inputs of the moment, and two
identical inputs give identical outputs, whatever happened before. A light, a \omterm{def:b1:sequential-circuits:register}{counter}, a lock and a processor all need
\emph{memory}. This chapter builds memory from gates by letting them feed back on themselves, tames it with a clock,
and arrives at the \omterm{def:b1:sequential-circuits:fsm}{finite-state machine}, the model behind every controller, whose most important example is the \omterm{def:b1:sequential-circuits:datapath}{control
unit} of a processor. The last section builds one for the toy machine and checks it, instruction by instruction, against
the toy machine's simulator.

\begin{recall}
The second book met the \omterm{def:b1:sequential-circuits:flip-flop}{flip-flop} and the state machine, and its toy machine runs a fetch--decode--execute cycle;
\omterm{def:b1:combinational-circuits:combinational}{combinational circuits}, \omterm{def:b1:combinational-circuits:karnaugh}{Karnaugh maps}, the ALU and \omterm{def:b1:combinational-circuits:delay}{propagation delay} are in \cref{ch:b1:combinational-circuits}.
\end{recall}

\section{Feedback and memory}\label{sec:b1:sequential-circuits:feedback}

\begin{definition}[Sequential circuit]\label{def:b1:sequential-circuits:sequential}
A \emph{sequential circuit}\index{sequential circuit} is a circuit whose outputs depend on the past inputs as well as the
present ones. Its gates form a \omterm{def:b1:graph-traversals:graph}{graph} with cycles, or it contains elements that store bits.
\end{definition}

The simplest has two NOR gates, each fed by the other. The simulator for such circuits cannot evaluate the gates in a
\omterm{def:b1:graph-traversals:topo}{topological order}, since there is none; it moves time forward one gate delay at a time, every gate taking the values its
inputs had one unit before.

\omcode{code/bachelor-1/28-sequential-circuits/listings/sr_latch.py}{7}{11}{Two NOR gates feeding each
other}\label{lst:b1:sequential-circuits:sr}

\omoutput{code/bachelor-1/28-sequential-circuits/outputs/sr_latch.txt}{The \omterm{def:b1:sequential-circuits:latch}{latch} traced one gate delay at a time.
After \emph{set} and \emph{reset} it holds its bit; released from $S = R = 1$ it
oscillates.}\label{out:b1:sequential-circuits:sr}

\begin{definition}[Latch, SR latch, D latch]\label{def:b1:sequential-circuits:latch}
A \emph{latch}\index{latch} is a circuit with two stable \omterm{def:b1:sequential-circuits:fsm}{states} that stores one bit. An \emph{SR latch}\index{SR latch} has
inputs $S$ (set) and $R$ (reset): $S = 1$ makes its output $Q$ 1, $R = 1$ makes it 0, and $S = R = 0$ keeps it. The
input $S = R = 1$ is forbidden. A \emph{D latch}\index{D latch} has a data input $D$ and an enable $E$: while $E = 1$ its
output follows $D$ (it is \emph{transparent}), and while $E = 0$ it keeps the last value.
\end{definition}

\begin{omfigure}
\begin{minipage}{0.55\linewidth}\centering
\begin{circuitikz}[font=\small]
\node[nor port] (top) at (2,1) {}; \node[nor port] (bot) at (2,-1) {};
\draw (top.in 1) -- ++(-0.9,0) node[left] {$R$}; \draw (bot.in 2) -- ++(-0.9,0) node[left] {$S$};
\draw (top.out) -- ++(0.9,0) node[right] {$Q$}; \draw (bot.out) -- ++(0.9,0) node[right] {$Q'$};
\fill ($(top.out)+(0.3,0)$) circle (1.3pt) ($(bot.out)+(0.3,0)$) circle (1.3pt);
\draw ($(top.out)+(0.3,0)$) -- ++(0,-0.6) -- ($(bot.in 1)+(-0.4,0.45)$) |- (bot.in 1);
\draw ($(bot.out)+(0.3,0)$) -- ++(0,0.6) -- ($(top.in 2)+(-0.4,-0.45)$) |- (top.in 2);
\end{circuitikz}
\end{minipage}%
\begin{minipage}{0.4\linewidth}\centering\small
\begin{tabular}{cc|l}
$S$ & $R$ & $Q$ next \\ \hline
0 & 0 & $Q$ (hold) \\
1 & 0 & 1 (set) \\
0 & 1 & 0 (reset) \\
1 & 1 & forbidden
\end{tabular}
\end{minipage}
\par\smallskip{\small An \omterm{def:b1:sequential-circuits:latch}{SR latch}: the output of each NOR gate is an input of the other. The loop is what remembers.}
\end{omfigure}

With $S = R = 0$ each NOR gate is an inverter, and two inverters in a loop are stable in two ways, $Q = 1, Q' = 0$ and
$Q = 0, Q' = 1$: that is the stored bit. $S = 1$ forces $Q' = 0$, hence $Q = 1$; $R = 1$ forces the opposite. With $S = R
= 1$ both outputs are 0, which is not a pair of complements; releasing both inputs at once lets the two gates race,
and in a model with exactly equal delays they oscillate for ever. In real gates one wins, but nobody can say which: that is
why the input is forbidden.

A \omterm{def:b1:sequential-circuits:latch}{D latch} puts two AND gates in front: $S = DE$ and $R = D'E$, so the forbidden input cannot happen. Its transparency is a
problem, however. If the output of a \omterm{def:b1:sequential-circuits:latch}{latch} feeds, through some logic, back to its own input, a change can run round the
loop many times while the \omterm{def:b1:sequential-circuits:latch}{latch} is enabled, and the result depends on how long the enable lasts. The cure is to let the
stored value change only at an instant.

\section{Flip-flops and the clock}\label{sec:b1:sequential-circuits:clock}

\begin{definition}[Flip-flop, edge-triggered D flip-flop]\label{def:b1:sequential-circuits:flip-flop}
A \emph{flip-flop}\index{flip-flop} stores one bit and changes it only at an edge of a control signal. An
\emph{edge-triggered D flip-flop}\index{edge-triggered D flip-flop} copies its input $D$ to its output $Q$ at each
rising edge (a change from 0 to 1) of its clock input, and keeps $Q$ the rest of the time.
\end{definition}

Two \omterm{def:b1:sequential-circuits:latch}{D latches} make one. The first, the \emph{master}, is enabled while the clock is 0 and follows $D$; the second, the
\emph{slave}, is enabled while the clock is 1 and copies the master. At the rising edge the master closes, holding the
value $D$ had at that moment, and the slave opens and shows it; while the clock is 1, $D$ may change, but the master is
closed.

\omcode{code/bachelor-1/28-sequential-circuits/listings/d_flipflop.py}{10}{25}{A master--slave D \omterm{def:b1:sequential-circuits:flip-flop}{flip-flop} from two \omterm{def:b1:sequential-circuits:latch}{D
latches}}\label{lst:b1:sequential-circuits:dff}

\omoutput{code/bachelor-1/28-sequential-circuits/outputs/d_flipflop.txt}{The \omterm{def:b1:sequential-circuits:flip-flop}{flip-flop} under a clock of period 20 time
units: $Q$ changes only after rising edges. Then $D$ is changed at $k$ units from an edge, to find the window in which it
must not move.}\label{out:b1:sequential-circuits:dff}

\begin{definition}[Clock signal, clock period, synchronous circuit]\label{def:b1:sequential-circuits:clock}
A \emph{clock signal}\index{clock signal} is a signal that alternates between 0 and 1 at a fixed rate; the time between
two rising edges is the \emph{clock period}\index{clock period}, and its inverse the clock frequency. A
\emph{synchronous circuit}\index{synchronous circuit} is made of \omterm{def:b1:sequential-circuits:flip-flop}{flip-flops}, all driven by the same clock, and of
combinational logic between them, with no loop that does not pass through a \omterm{def:b1:sequential-circuits:flip-flop}{flip-flop}.
\end{definition}

In a \omterm{def:b1:sequential-circuits:clock}{synchronous circuit} the state, the contents of the \omterm{def:b1:sequential-circuits:flip-flop}{flip-flops}, changes only at the edges; between two edges the
combinational logic computes, from the state and the inputs, the next state, which the \omterm{def:b1:sequential-circuits:flip-flop}{flip-flops} load at the next edge.
Every loop is cut by a \omterm{def:b1:sequential-circuits:flip-flop}{flip-flop}, so the logic between edges is combinational, and \cref{ch:b1:combinational-circuits}
applies to it.

\begin{definition}[Setup time, hold time]\label{def:b1:sequential-circuits:setup}
The \emph{setup time}\index{setup time} of a \omterm{def:b1:sequential-circuits:flip-flop}{flip-flop} is how long before the clock edge its input must be stable; the
\emph{hold time}\index{hold time}, how long after the edge it must stay stable. The \emph{clock-to-output delay} is the
time from the edge to the new output.
\end{definition}

The measurement in \cref{out:b1:sequential-circuits:dff} finds a \omterm{def:b1:sequential-circuits:setup}{setup time} of 2 units: a $D$ that falls one unit before
the edge is missed, because the path from $D$ through the NOT gate into the master is one gate longer than the path of a
rising $D$. Its \omterm{def:b1:sequential-circuits:setup}{hold time} is 0, and the clock-to-output delay 2 units. The figure below draws the run.

\begin{omfigure}
\begin{tikzpicture}
\begin{axis}[width=\linewidth, height=46mm, xmin=0, xmax=79, ymin=-0.3, ymax=5.3, axis lines=none, clip=false]
\addplot[const plot, thick, black] table[x=t, y expr=\thisrow{clk}*0.9+4]
  {code/bachelor-1/28-sequential-circuits/figures/dff.dat};
\addplot[const plot, thick, omDef] table[x=t, y expr=\thisrow{D}*0.9+2]
  {code/bachelor-1/28-sequential-circuits/figures/dff.dat};
\addplot[const plot, thick, omThm] table[x=t, y expr=\thisrow{Q}*0.9]
  {code/bachelor-1/28-sequential-circuits/figures/dff.dat};
\node[left] at (axis cs:0,4.45) {clock}; \node[left] at (axis cs:0,2.45) {$D$}; \node[left] at (axis cs:0,0.45) {$Q$};
\pgfplotsinvokeforeach{10,30,50,70}{
  \draw[gray, dashed] (axis cs:#1,-0.2) -- (axis cs:#1,5.1);
  \fill[omThm!25, opacity=0.6] (axis cs:#1-2,1.9) rectangle (axis cs:#1,3.0);}
\node[font=\scriptsize, anchor=south] at (axis cs:50,5.0) {rising edges (dashed)};
\node[font=\scriptsize, fill=white, inner sep=1pt] at (axis cs:60,3.45) {this pulse of $D$ is never stored};
\end{axis}
\end{tikzpicture}
\par\smallskip{\small The run of \cref{out:b1:sequential-circuits:dff}. $Q$ takes the value $D$ had at each rising edge,
two time units later; the shaded windows are the \omterm{def:b1:sequential-circuits:setup}{setup times}, in which $D$ must not change.}
\end{omfigure}

The \omterm{def:b1:sequential-circuits:clock}{clock period} is bounded below. After an edge, the new state appears at the \omterm{def:b1:sequential-circuits:flip-flop}{flip-flop} outputs a clock-to-output delay
later, crosses the combinational logic along its \omterm{def:b1:combinational-circuits:delay}{longest delay path} (\cref{def:b1:combinational-circuits:delay}), and
must arrive a \omterm{def:b1:sequential-circuits:setup}{setup time} before the next edge:
\[ T \geq t_{\text{clock-to-output}} + t_{\text{logic}} + t_{\text{setup}}. \]
The \omterm{def:b1:sequential-circuits:setup}{hold time} puts a bound on the \emph{shortest} path: the new values must not reach the next \omterm{def:b1:sequential-circuits:flip-flop}{flip-flops} before their \omterm{def:b1:sequential-circuits:setup}{hold
time} has passed. A processor's clock frequency is set by the slowest path between two of its \omterm{def:b1:sequential-circuits:register}{registers}.

\section{Registers, counters and memory}\label{sec:b1:sequential-circuits:registers}

\begin{definition}[Register, shift register, counter]\label{def:b1:sequential-circuits:register}
A \emph{register}\index{register (hardware)} is a row of \omterm{def:b1:sequential-circuits:flip-flop}{flip-flops} on the same clock that store a word; with an
\emph{enable} input, each loads its input when enable is 1 and keeps its value otherwise. A \emph{shift
register}\index{shift register} moves its bits one place at each edge, a new bit entering at one end. A
\emph{counter}\index{counter (circuit)} is a register whose value goes through $0, 1, 2, \dots$ at successive edges.
\end{definition}

The enable is a \omterm{def:b1:combinational-circuits:mux}{multiplexer} in front of each \omterm{def:b1:sequential-circuits:flip-flop}{flip-flop}, choosing between its own output and the new input: a \omterm{def:b1:sequential-circuits:register}{register}
with enable is a \omterm{def:b1:sequential-circuits:clock}{synchronous circuit}, its clock never stopped.

\omcode{code/bachelor-1/28-sequential-circuits/listings/shift.py}{6}{17}{A \omterm{def:b1:sequential-circuits:register}{register} with enable and a \omterm{def:b1:sequential-circuits:register}{shift register}, one
clock edge at a time}\label{lst:b1:sequential-circuits:shift}

\omoutput{code/bachelor-1/28-sequential-circuits/outputs/shift.txt}{Loading only when enabled; four bits entering one by
one, then moving on.}\label{out:b1:sequential-circuits:shift}

A \omterm{def:b1:sequential-circuits:register}{shift register} turns a stream of bits into a word (after four edges the bits 1, 1, 0, 1 sit side by side) and back:
that is how a serial line, such as a USB cable, carries bytes one bit at a time.

A \omterm{def:b1:sequential-circuits:register}{counter} can be built two ways. In a \emph{synchronous \omterm{def:b1:sequential-circuits:register}{counter}} every \omterm{def:b1:sequential-circuits:flip-flop}{flip-flop} is on the clock, and gates compute the
next value: bit 0 flips at every edge, bit 1 when bit 0 is 1, bit 2 when bits 0 and 1 are both 1. In a \emph{ripple
\omterm{def:b1:sequential-circuits:register}{counter}} only the first \omterm{def:b1:sequential-circuits:flip-flop}{flip-flop} sees the clock; each of the others flips when the one before it falls from 1 to 0. It
needs no gates, but its bits change one after the other.

\omcode{code/bachelor-1/28-sequential-circuits/listings/counter.py}{6}{33}{Two 3-bit \omterm{def:b1:sequential-circuits:register}{counters}, watched at every time
unit; each \omterm{def:b1:sequential-circuits:flip-flop}{flip-flop} takes $d$ units}\label{lst:b1:sequential-circuits:counter}

\omoutput{code/bachelor-1/28-sequential-circuits/outputs/counter.txt}{The values seen between two clock edges. The ripple
\omterm{def:b1:sequential-circuits:register}{counter} passes through wrong values: from 3 to 4 it shows 2 and 0 on the
way.}\label{out:b1:sequential-circuits:counter}

A circuit that reads the ripple \omterm{def:b1:sequential-circuits:register}{counter} while it settles may see 2 or 0 between 3 and 4: a \emph{glitch}. The synchronous
\omterm{def:b1:sequential-circuits:register}{counter} has none, and settles in one \omterm{def:b1:sequential-circuits:flip-flop}{flip-flop} delay whatever its width, where the ripple \omterm{def:b1:sequential-circuits:register}{counter} needs $n$.

A memory is many \omterm{def:b1:sequential-circuits:register}{registers} and a way to choose one. A $k$-to-$2^k$ \omterm{def:b1:combinational-circuits:mux}{decoder} (\cref{def:b1:combinational-circuits:mux})
turns the address into one selected row; a write loads the row whose \omterm{def:b1:combinational-circuits:mux}{decoder} output is 1; a read sends the selected row
through a \omterm{def:b1:combinational-circuits:mux}{multiplexer} to the output.

\begin{definition}[Memory cell]\label{def:b1:sequential-circuits:memory-cell}
A \emph{memory cell}\index{memory cell} is the circuit that stores one bit of a memory, with an input that selects it for
writing and an output that the read logic selects.
\end{definition}

\omcode{code/bachelor-1/28-sequential-circuits/listings/ram.py}{7}{26}{The \omterm{def:b1:combinational-circuits:mux}{decoder} and the \omterm{def:b1:combinational-circuits:mux}{multiplexer} of a 16-word
memory}\label{lst:b1:sequential-circuits:ram}

\omoutput{code/bachelor-1/28-sequential-circuits/outputs/ram.txt}{A $16 \times 4$ memory: writes and reads, and its
gates.}\label{out:b1:sequential-circuits:ram}

More than half of the gates are cells, and the share grows with the size: a memory of $2^k$ words of $w$ bits needs $5
\cdot 2^k w$ gates of cells against about $2^k$ for the \omterm{def:b1:combinational-circuits:mux}{decoder}. Real memories make the cell as small as possible: six
transistors in a static \omterm{def:b1:sequential-circuits:memory-cell}{memory cell}, one transistor and a capacitor, refreshed every few hundredths of a second, in a dynamic
one.

\begin{history}[Magnetic cores]
% ledger: b1:core-memory
Before memories were made of transistors, the main memory of computers was a plane of tiny rings of magnetic material,
threaded on a grid of wires: each ring held one bit as the direction of its magnetisation, and the current in one row
wire and one column wire together selected a single ring, a \omterm{def:b1:combinational-circuits:mux}{decoder} made of physics. Jay Forrester described the
three-dimensional arrangement in 1951, and core memory was the main memory of most computers until the 1970s.
\end{history}

\omimg{book3/core-memory-plane.jpg}{0.4\linewidth}{A plane of core memory holding 1024 bits: 32 rows and 32 columns of
ferrite rings, each threaded by the wires of its row and its column. Photograph: Konstantin Lanzet, CC BY-SA 3.0.}

\section{Finite-state machines}\label{sec:b1:sequential-circuits:fsm}

Every \omterm{def:b1:sequential-circuits:clock}{synchronous circuit} is the same object seen from far: a state, kept in \omterm{def:b1:sequential-circuits:flip-flop}{flip-flops}; an input at each clock cycle;
a next state and an output computed from them. Mathematics names it.

\begin{definition}[Finite-state machine]\label{def:b1:sequential-circuits:fsm}
A \emph{finite-state machine}\index{finite-state machine} has a finite set of \emph{states}\index{state (of a machine)},
one of them initial, a finite input alphabet $\Sigma$, an output alphabet, and a \emph{transition}\index{transition}
function that gives, for each state $q$ and input symbol $a$, the next state $\delta(q, a)$. It reads one symbol per step
and produces one output per step.
\end{definition}

\begin{definition}[Moore machine, Mealy machine, state diagram, state encoding]\label{def:b1:sequential-circuits:moore-mealy}
In a \emph{Moore machine}\index{Moore machine} the output is a function $\lambda(q)$ of the state; in a \emph{Mealy
machine}\index{Mealy machine} it is a function $\lambda(q, a)$ of the state and the input, attached to the \omterm{def:b1:sequential-circuits:fsm}{transition}. A
\emph{state diagram}\index{state diagram} draws the \omterm{def:b1:sequential-circuits:fsm}{states} as circles and each \omterm{def:b1:sequential-circuits:fsm}{transition} as an arrow labelled by its
input (and, for a Mealy machine, its output, as $a/o$). A \emph{state encoding}\index{state encoding} gives each state a
distinct word of bits, to be stored in \omterm{def:b1:sequential-circuits:flip-flop}{flip-flops}.
\end{definition}

For an input word $a_1 \dots a_n$ a \omterm{def:b1:sequential-circuits:moore-mealy}{Mealy machine} outputs $\lambda(q_0, a_1) \dots \lambda(q_{n-1}, a_n)$ and a \omterm{def:b1:sequential-circuits:moore-mealy}{Moore
machine} $\lambda(q_1) \dots \lambda(q_n)$, where $q_i$ is the state after $i$ symbols. Two machines are
\emph{equivalent} if they give the same output word for every input word. Both kinds compute the same things:

\begin{proposition}[Moore and Mealy are equivalent]\label{prop:b1:sequential-circuits:conversion}
For every \omterm{def:b1:sequential-circuits:moore-mealy}{Moore machine} there is an equivalent \omterm{def:b1:sequential-circuits:moore-mealy}{Mealy machine} with the same \omterm{def:b1:sequential-circuits:fsm}{states}, and for every \omterm{def:b1:sequential-circuits:moore-mealy}{Mealy machine} with state
set $Q$ and output alphabet $O$ there is an equivalent \omterm{def:b1:sequential-circuits:moore-mealy}{Moore machine} with at most $|Q| \cdot |O|$ \omterm{def:b1:sequential-circuits:fsm}{states}.
\end{proposition}

\begin{proof}
Moore to Mealy: put on each \omterm{def:b1:sequential-circuits:fsm}{transition} the output of the state it enters, $\lambda'(q, a) = \lambda(\delta(q, a))$. Mealy to
Moore: take the pairs $(q, o)$ as \omterm{def:b1:sequential-circuits:fsm}{states}, meaning ``in $q$, having just output $o$'', with $\delta'((q, o), a) = (\delta(q,
a), \lambda(q, a))$ and output $o$. Both machines go through the same \omterm{def:b1:sequential-circuits:fsm}{states} $q_i$, and the output read after step $i$ is
$\lambda(q_{i-1}, a_i)$ in both; \cref{exo:b1:sequential-circuits:10} writes the induction out.
\end{proof}

\omcode{code/bachelor-1/28-sequential-circuits/listings/moore_mealy.py}{35}{52}{The two conversions, keeping only the
reachable pairs}\label{lst:b1:sequential-circuits:conversion}

The example is a \emph{sequence detector}: it reads bits and outputs 1 exactly when the last three were $1, 0, 1$.
Overlaps count: in $10101$ it fires twice. Its Mealy \omterm{def:b1:sequential-circuits:fsm}{states} remember the longest end of the input that is a beginning of
$101$: nothing, $1$, or $10$.

\omoutput{code/bachelor-1/28-sequential-circuits/outputs/moore_mealy.txt}{The \omterm{def:b1:sequential-circuits:moore-mealy}{Moore machine} made from the Mealy detector of
101: the state $(1, 1)$ is new, the state $1$ entered with an output of 1. Both agree on every word of up to ten
bits.}\label{out:b1:sequential-circuits:conversion}

\begin{omfigure}
\begin{tikzpicture}[->, >=stealth, auto, node distance=21mm, every state/.style={minimum size=8mm, font=\small},
  every edge quotes/.style={font=\scriptsize}]
\node[state, initial, initial text={}] (e) {$\varepsilon$}; \node[state, right=of e] (a) {$1$};
\node[state, right=of a] (b) {$10$};
\path (e) edge[loop above] node {$0/0$} ()
      (e) edge node {$1/0$} (a)
      (a) edge[loop above] node {$1/0$} ()
      (a) edge[bend left=20] node {$0/0$} (b)
      (b) edge[bend left=20] node {$1/\mathbf{1}$} (a)
      (b) edge[bend left=40] node {$0/0$} (e);
\begin{scope}[xshift=78mm]
\node[state, initial, initial text={}] (E) {$\varepsilon$/0}; \node[state, right=of E] (A) {$1$/0};
\node[state, right=of A] (B) {$10$/0}; \node[state, accepting, below=14mm of A, fill=omThm!15] (C) {$1$/1};
\path (E) edge[loop above] node {0} () (E) edge node {1} (A) (A) edge[loop above] node {1} ()
      (A) edge node {0} (B) (B) edge node {1} (C) (B) edge[bend right=62, swap] node {0} (E)
      (C) edge node {0} (B) (C) edge node {1} (A);
\end{scope}
\end{tikzpicture}
\par\smallskip{\small The detector of 101. Left, the \omterm{def:b1:sequential-circuits:moore-mealy}{Mealy machine}: the 1 is output on the \omterm{def:b1:sequential-circuits:fsm}{transition} from $10$ on input 1.
Right, the \omterm{def:b1:sequential-circuits:moore-mealy}{Moore machine}: the 1 is the output of the state $1/1$, which splits the state $1$ in two.}
\end{omfigure}

\paragraph{From a \omterm{def:b1:sequential-circuits:moore-mealy}{state diagram} to a circuit.} Four steps make a machine into a \omterm{def:b1:sequential-circuits:clock}{synchronous circuit}. Encode the \omterm{def:b1:sequential-circuits:fsm}{states} in
$\lceil \log_2 |Q| \rceil$ bits, one \omterm{def:b1:sequential-circuits:flip-flop}{flip-flop} per bit; write the truth tables of the next-state bits and of the output,
as functions of the state bits and the input; simplify them (\cref{sec:b1:combinational-circuits:simplification}); and
connect the logic's outputs to the \omterm{def:b1:sequential-circuits:flip-flop}{flip-flops}' inputs. Unused codes are don't-cares.

\omoutput{code/bachelor-1/28-sequential-circuits/outputs/detector.txt}{The Moore detector with $(\varepsilon, 0) = 00$,
$(1, 0) = 01$, $(10, 0) = 10$, $(1, 1) = 11$; $A$ is the high state bit, $B$ the low one, $x$ the input. The circuit is
checked against the machine.}\label{out:b1:sequential-circuits:detector}

$B^+ = x$ says that the low state bit is just the last input, a one-bit \omterm{def:b1:sequential-circuits:register}{shift register}; $z = AB$ needs one gate. A
different encoding gives different equations, sometimes better ones; \emph{one-hot} encoding, one \omterm{def:b1:sequential-circuits:flip-flop}{flip-flop} per state,
spends \omterm{def:b1:sequential-circuits:flip-flop}{flip-flops} to make the logic trivial (\cref{exo:b1:sequential-circuits:9}).

\section{The control of a small processor}\label{sec:b1:sequential-circuits:control}

The toy machine of the earlier books executes an instruction in one step of its simulator. Hardware cannot: reading the
two bytes of the instruction, reading or writing memory, and changing \omterm{def:b1:sequential-circuits:register}{registers} take several clock cycles, and something
must order them. A processor is two parts.

\begin{definition}[Datapath, control unit]\label{def:b1:sequential-circuits:datapath}
The \emph{datapath}\index{datapath} of a processor is the part that holds and transforms data: \omterm{def:b1:sequential-circuits:register}{registers}, the ALU,
the buses and \omterm{def:b1:combinational-circuits:mux}{multiplexers} between them, and the connection to memory. The \emph{control unit}\index{control unit} is a
\omterm{def:b1:sequential-circuits:fsm}{finite-state machine} that, at each clock cycle, sends the datapath the signals that choose which \omterm{def:b1:sequential-circuits:register}{register} transfer
happens: which \omterm{def:b1:sequential-circuits:register}{register} is loaded, from which source, with which ALU operation.
\end{definition}

The \omterm{def:b1:sequential-circuits:datapath}{datapath} of the toy machine has the four \omterm{def:b1:sequential-circuits:register}{registers}, the program \omterm{def:b1:sequential-circuits:register}{counter} \code{PC}, the stack \omterm{def:b1:pointers-and-arrays:pointer}{pointer} \code{SP}, the
flags, and two instruction \omterm{def:b1:sequential-circuits:register}{registers} \code{IR0} and \code{IR1} for the two bytes of the current instruction; its ALU is
that of \cref{def:b1:combinational-circuits:alu} on eight bits, which the program computes with integer arithmetic,
except that, as in the toy machine, the flag C after a subtraction is the borrow. The
\omterm{def:b1:sequential-circuits:datapath}{control unit} is a \omterm{def:b1:sequential-circuits:moore-mealy}{Moore machine}: each state names the \omterm{def:b1:sequential-circuits:register}{register} transfers it orders. Its next state depends on the state,
on the operation code in \code{IR0}, and, for a conditional jump, on the flags.

\omcode{code/bachelor-1/28-sequential-circuits/listings/toy_control.py}{10}{41}{The \omterm{def:b1:sequential-circuits:datapath}{control unit} of the toy machine:
the transfers each state orders, and the next-state function}\label{lst:b1:sequential-circuits:control}

\code{R} is the \omterm{def:b1:sequential-circuits:register}{register} named in the first byte, \code{S} the one named in the second, and \code{M[...]} the memory.
\code{FETCH0} and \code{FETCH1} read the two bytes of the instruction; \code{DECODE} chooses the execute state from the
operation code; each execute state goes back to \code{FETCH0}. A jump not taken goes straight back too.

\omcode{code/bachelor-1/28-sequential-circuits/listings/toy_control.py}{115}{132}{Clocking the two parts, and comparing
the \omterm{def:b1:sequential-circuits:datapath}{datapath} with the simulator after every instruction}\label{lst:b1:sequential-circuits:clocking}

\omoutput{code/bachelor-1/28-sequential-circuits/outputs/toy_control.txt}{The programs of the earlier books, run by the
\omterm{def:b1:sequential-circuits:datapath}{control unit} and its \omterm{def:b1:sequential-circuits:datapath}{datapath}; after every instruction the \omterm{def:b1:sequential-circuits:register}{registers}, flags, memory and outputs are those of the toy
machine's simulator.}\label{out:b1:sequential-circuits:control}

\begin{omfigure}
\begin{tikzpicture}[->, >=stealth, st/.style={draw, rounded corners=2mm, font=\scriptsize\ttfamily, minimum height=6mm,
  inner sep=2pt}, font=\scriptsize]
\node[st, fill=omDef!12] (f0) at (0,0) {FETCH0}; \node[st, fill=omDef!12] (f1) at (2.4,0) {FETCH1};
\node[st, fill=omDef!12] (dec) at (4.8,0) {DECODE};
\foreach \n/\l/\y in {x1/LOADI/2.1, x2/LOAD/1.5, x3/STORE/0.9, x4/ALU\_RR/0.3, x5/CMP/-0.3, x6/JUMP/-0.9,
  x7/OUT/-1.5, x8/RET/-2.1}
  {\node[st] (\n) at (8.2,\y) {\l}; \draw (dec.east) -- (\n.west);}
\node[st] (p1) at (10.6,-0.3) {PUSH1}; \node[st] (p2) at (12.4,-0.3) {PUSH2};
\draw (dec.south) |- (10.6,-2.75) -- (p1.south); \draw (p1) -- (p2);
\node[st] (h) at (10.6,1.5) {HALT}; \node[st] (s) at (12.4,1.5) {STOP};
\draw (dec.north) |- (10.6,2.75) -- (h.north); \draw (h) -- (s);
\draw (f0) -- (f1); \draw (f1) -- (dec);
\draw (dec.south west) to[bend left=35] node[below] {jump not taken} (f0.south east);
\draw[gray] (8.2,-2.45) -- (8.2,-3.2) -- node[below] {every execute state returns to FETCH0} (0,-3.2) -- (f0.south west);
\node[anchor=west, align=left, gray] at (8.9,0.3) {\dots\ and ALU\_RI, ALU\_R, IN,\\ LOADR, STORER, POP, CALL1--2};
\end{tikzpicture}
\par\smallskip{\small The \omterm{def:b1:sequential-circuits:datapath}{control unit} of \cref{lst:b1:sequential-circuits:control}, some execute \omterm{def:b1:sequential-circuits:fsm}{states} left out: four
\omterm{def:b1:sequential-circuits:fsm}{states} per instruction for most (two fetch \omterm{def:b1:sequential-circuits:fsm}{states}, decode, one execute state), three for a jump not
taken, five for PUSH and CALL. The program counts 3.9 cycles per instruction
on average.}
\end{omfigure}

The \omterm{def:b1:sequential-circuits:datapath}{control unit} has 23 \omterm{def:b1:sequential-circuits:fsm}{states}, so five \omterm{def:b1:sequential-circuits:flip-flop}{flip-flops} encode it, and its next-state and output functions are combinational
logic of the state, six bits of operation code and three flags, which \cref{sec:b1:sequential-circuits:fsm} turns into
gates. Large instruction sets made those functions unmanageable; the answer, \emph{microprogrammed control},
stores the control signals of each state as a word in a small read-only memory, the \emph{microprogram}, addressed by the
state, so that the \omterm{def:b1:sequential-circuits:datapath}{control unit} is itself a tiny computer reading a table. Many processors since have been built that
way, and many have gone back to gates when speed mattered more than the size of the instruction set.

\section{Exercises}\label{sec:b1:sequential-circuits:exercises}

\begin{exercise}[$\star$]\label{exo:b1:sequential-circuits:1}
An \omterm{def:b1:sequential-circuits:latch}{SR latch} starts with $Q = 0$. Its inputs $(S, R)$ take the values $10$, $00$, $01$, $10$, $00$, each held long enough
for the \omterm{def:b1:sequential-circuits:latch}{latch} to settle. Give $Q$ after each.
\end{exercise}

\begin{exercise}[$\star$]\label{exo:b1:sequential-circuits:2}
The synchronous 3-bit \omterm{def:b1:sequential-circuits:register}{counter} of \cref{lst:b1:sequential-circuits:counter} holds 5. Give its value after each of the next
five clock edges, and which bits flip at each.
\end{exercise}

\begin{exercise}[$\star$]\label{exo:b1:sequential-circuits:3}
What does the Moore detector of \cref{out:b1:sequential-circuits:conversion} output on the input word $11010110$? Give the
\omterm{def:b1:sequential-circuits:fsm}{states} it goes through.
\end{exercise}

\begin{exercise}[$\star$]\label{exo:b1:sequential-circuits:4}
A 4-bit \omterm{def:b1:sequential-circuits:register}{shift register} as in \cref{lst:b1:sequential-circuits:shift} holds $1010$. The bits $1, 1, 0, 0$ enter at four
edges. Give its contents after each.
\end{exercise}

\begin{exercise}[$\star\star$]\label{exo:b1:sequential-circuits:5}
A circuit compares the 3-bit ripple \omterm{def:b1:sequential-circuits:register}{counter} of \cref{out:b1:sequential-circuits:counter} with 2, and sounds a bell when
they are equal. The \omterm{def:b1:sequential-circuits:register}{counter} goes from 0 to 7 once. How many times does the bell sound, and when? Repair the design.
\end{exercise}

\begin{exercise}[$\star\star$]\label{exo:b1:sequential-circuits:6}
This traffic light stops with an error on the second input word. Explain why, and complete the machine.
\omcode{code/bachelor-1/28-sequential-circuits/listings/exo_incomplete.py}{4}{15}{A traffic light}
\omoutput{code/bachelor-1/28-sequential-circuits/outputs/exo_incomplete.txt}{The two runs.}
\end{exercise}

\begin{exercise}[$\star\star$]\label{exo:b1:sequential-circuits:7}
Build the Mealy detector of $0110$, overlaps allowed, and test it on every word of up to 12 bits against the definition.
\end{exercise}

\begin{exercise}[$\star\star$]\label{exo:b1:sequential-circuits:8}
Design a synchronous \omterm{def:b1:sequential-circuits:register}{counter} modulo 6 on three D \omterm{def:b1:sequential-circuits:flip-flop}{flip-flops}: the next-state equations, simplified with the codes 6 and 7
as don't-cares, and a check that it counts $0, 1, \dots, 5, 0, \dots$.
\end{exercise}

\begin{exercise}[$\star\star$]\label{exo:b1:sequential-circuits:9}
How many \omterm{def:b1:sequential-circuits:flip-flop}{flip-flops} does a machine with $k$ \omterm{def:b1:sequential-circuits:fsm}{states} need with a binary encoding? With one-hot encoding? For the \omterm{def:b1:sequential-circuits:datapath}{control
unit} of \cref{lst:b1:sequential-circuits:control}, compare, and say what one-hot encoding buys.
\end{exercise}

\begin{exercise}[$\star\star\star$]\label{exo:b1:sequential-circuits:10}
Prove by induction on the length of the input word that the conversions of
\cref{prop:b1:sequential-circuits:conversion} give equivalent machines.
\end{exercise}

\begin{exercise}[$\star\star\star$]\label{exo:b1:sequential-circuits:11}
In a \omterm{def:b1:sequential-circuits:clock}{synchronous circuit} the \omterm{def:b1:sequential-circuits:flip-flop}{flip-flops} have a clock-to-output delay of 50 ps, a \omterm{def:b1:sequential-circuits:setup}{setup time} of 30 ps and a \omterm{def:b1:sequential-circuits:setup}{hold time} of
10 ps; every gate takes 25 ps. Between two \omterm{def:b1:sequential-circuits:register}{registers} the longest path has 12 gates and the shortest 2. What is the
highest clock frequency? Is the \omterm{def:b1:sequential-circuits:setup}{hold time} respected? What if the \omterm{def:b1:shortest-paths:shortest}{shortest path} had no gate?
\end{exercise}

\begin{exercise}[$\star\star\star$]\label{exo:b1:sequential-circuits:12}
A vending machine sells an item for 3; at each clock cycle a coin of 1 or 2, or none, is inserted. Draw its \omterm{def:b1:sequential-circuits:moore-mealy}{Mealy
machine} (outputs: give the item, give back 1), derive simplified next-state and output equations, and simulate the circuit
against the machine on every sequence of up to eight coins.
\end{exercise}

\section{Problem: A combination lock}\label{sec:b1:sequential-circuits:problem}

\begin{problem}[{Weekend problem --- A combination lock}]\label{pb:b1:sequential-circuits:1}
A lock has two keys, 0 and 1, and opens when the last four keys pressed are $1, 0, 1, 1$. The program builds its \omterm{def:b1:sequential-circuits:moore-mealy}{Mealy
machine} from the pattern, converts it to a \omterm{def:b1:sequential-circuits:moore-mealy}{Moore machine}, finds for each pair of \omterm{def:b1:sequential-circuits:fsm}{states} a word that tells them apart,
and turns both machines into circuits, checked on every word of up to 12 keys.
\omcode{code/bachelor-1/28-sequential-circuits/listings/lock.py}{13}{24}{The detector of any pattern: \omterm{def:b1:sequential-circuits:fsm}{states} are the
prefixes of the pattern}\label{lst:b1:sequential-circuits:lock}
\omoutput{code/bachelor-1/28-sequential-circuits/outputs/lock.txt}{The two machines, the words that distinguish their
\omterm{def:b1:sequential-circuits:fsm}{states} ($e$ is the empty prefix; in the \omterm{def:b1:sequential-circuits:moore-mealy}{Moore machine} $p/o$ is the prefix $p$ entered with output $o$), and their
circuits.}\label{out:b1:sequential-circuits:lock}

\needspace{4\baselineskip}
\medskip\noindent\textbf{Part I --- The Mealy machine.}
\begin{enumerate}
\item Draw the \omterm{def:b1:sequential-circuits:moore-mealy}{Mealy machine} of the output: four \omterm{def:b1:sequential-circuits:fsm}{states}, eight \omterm{def:b1:sequential-circuits:fsm}{transitions}.
\item Why is the next state of $101$ on key 1 the state $1$, and not $e$?
\item Which line of the listing finds that next state, and how?
\item Run the machine by hand on $1011011$. When does the lock open?
\item Why is the machine \emph{complete}: why does every state have a \omterm{def:b1:sequential-circuits:fsm}{transition} for both keys?
\item The output is 1 on a \omterm{def:b1:sequential-circuits:fsm}{transition}. When, exactly, should the bolt move?
\end{enumerate}

\needspace{4\baselineskip}
\medskip\noindent\textbf{Part II --- The Moore machine.}
\begin{enumerate}[resume]
\item Draw the \omterm{def:b1:sequential-circuits:moore-mealy}{Moore machine}; which state is new?
\item Why does the conversion give five \omterm{def:b1:sequential-circuits:fsm}{states} and not $4 \times 2 = 8$?
\item On the same word $1011011$, compare the outputs of the two machines step by step.
\item Why do the pairs involving $1/1$ need only the empty word?
\item Explain why a word that tells two \omterm{def:b1:sequential-circuits:fsm}{states} apart proves that no equivalent machine can merge them.
\item Conclude that four Mealy \omterm{def:b1:sequential-circuits:fsm}{states} and five Moore \omterm{def:b1:sequential-circuits:fsm}{states} are each the fewest possible.
\end{enumerate}

\needspace{4\baselineskip}
\medskip\noindent\textbf{Part III --- Circuits.}
\begin{enumerate}[resume]
\item Give the binary encodings the program uses, in the order of the output.
\item Check the Mealy equation $B^+ = x$ on two \omterm{def:b1:sequential-circuits:fsm}{transitions}.
\item Why may the Moore equations use the codes 5, 6 and 7 freely?
\item Why is the Moore output $z = A$, with no gate at all?
\item Count the gates of each circuit from its equations, as the program does.
\item Which circuit is faster from key to output, and why does the answer not depend only on gates?
\end{enumerate}

\needspace{4\baselineskip}
\medskip\noindent\textbf{Part IV --- Checking.}
\begin{enumerate}[resume]
\item How many words of up to 12 keys are there?
\item Show that if a circuit and a machine have the same output and next state on every reachable state and key, they
agree on every word, of any length.
\item Every state is reachable in at most four keys. Why do the words of up to 12 keys check every state and key, and
more?
\item What would a wrong encoding (two \omterm{def:b1:sequential-circuits:fsm}{states} with the same code) do to the check?
\item A real keypad has ten keys. How does the machine change, and how many words of up to 12 keys are there now?
\item State the named result, \emph{Mealy saves a state}, with the counts of \omterm{def:b1:sequential-circuits:fsm}{states}, \omterm{def:b1:sequential-circuits:flip-flop}{flip-flops} and gates.
\end{enumerate}
\end{problem}
